\section{System-by-System Analysis}
\label{sec:systems}

We organize systems into six categories: production systems
(\S\ref{sec:production}), standalone memory servers
(\S\ref{sec:standalone}), framework-level memory
(\S\ref{sec:framework}), coding agent memory (\S\ref{sec:coding-agents}),
file-based memory (\S\ref{sec:filebased}), and research architectures
(\S\ref{sec:research}). For each system we identify the memory
types supported (mapped to the CoALA taxonomy~\cite{sumers2024coala}),
the retrieval architecture (mapped to the RAG
taxonomy~\cite{gao2023rag}), the lifecycle operations, and the
infrastructure requirements. Quantitative cross-comparisons appear in
\S\ref{sec:comparison}; here we focus on what makes each system
architecturally distinct.

%% ===============================================================
\subsection{Production Systems}
\label{sec:production}

\paragraph{Mem0.}
Mem0~\cite{choudhary2025mem0} implements a tiered memory architecture
backed by vector databases (Qdrant, Pinecone, or Chroma), a relational
store for metadata, and an optional Neo4j graph layer (the Mem0g
variant). Memory formation is LLM-dependent: every incoming
interaction is routed through an LLM that extracts facts, user
preferences, and personality traits, then decides whether to create,
update, or deduplicate entries. Retrieval combines vector similarity
with graph traversal and contextual priority scoring, placing Mem0 in
the Advanced RAG category. Forgetting follows a single-curve
exponential decay of low-relevance entries, with periodic consolidation
between short-term and long-term tiers. The authors report a 26\%
accuracy improvement over OpenAI's built-in memory on the LOCOMO
benchmark and 91\% lower p95 latency versus full-context approaches.

The key architectural tradeoff in Mem0 is its heavy infrastructure
requirement (2--3 external services) and its reliance on an LLM for
every memory write. This makes Mem0 well-suited for cloud-deployed,
multi-user SaaS products with dedicated operations teams, but less
appropriate for privacy-first or offline use cases. The Mem0g graph
variant adds entity-relationship traversal, but the graph is
static---entities and relationships exist only when the LLM explicitly
creates them, with no dynamic restructuring through use.

\paragraph{Zep / Graphiti.}
Graphiti~\cite{dempsey2025graphiti}, the engine underlying Zep's memory
service, constructs a hierarchical temporal knowledge graph on Neo4j
with three subgraph layers: community (cluster summaries), entity
(semantic entities and relationships), and episodic (raw conversation
episodes). Retrieval uses a triple-hybrid pipeline---cosine similarity,
Okapi BM25, and breadth-first graph traversal---fused via Reciprocal
Rank Fusion (RRF) with subsequent maximal marginal relevance (MMR)
diversification. Multi-stage reranking considers episode mentions, node
distance, and cross-signal fusion. This places Zep/Graphiti firmly in
the Modular RAG category, alongside a small minority of systems
(17.9\%, Table~\ref{tab:rag-distribution}).

Zep's most distinctive lifecycle feature is \emph{bi-temporal
invalidation}: facts are never deleted, only marked with validity
periods recording when each fact was true in the real world and when
the system learned it. This conservative approach preserves full
provenance but means the knowledge graph grows monotonically. Entity
extraction requires an LLM (typically \texttt{gpt-4o-mini}), though
retrieval itself is LLM-free. The authors report 94.8\% accuracy on
the DMR benchmark and up to 18.5\% accuracy improvement on
LongMemEval.

\paragraph{Letta (formerly MemGPT).}
Letta~\cite{packer2023memgpt} introduced the operating-system metaphor
for agent memory: core memory serves as ``RAM'' (always in-context,
editable by the agent), recall memory stores conversation history, and
archival memory provides unbounded ``disk'' storage via a vector or
graph database. The defining architectural choice is that the LLM
\emph{is} the memory manager---it decides what to store, evict, and
retrieve via explicit tool calls (\texttt{core\_memory\_append},
\texttt{archival\_memory\_search}, etc.).

This agent-driven design makes Letta's memory quality entirely
dependent on LLM capability. When the context window fills, the
system evicts approximately 70\% of messages through summarization.
A 2025 extension adds ``sleep-time agents'' that handle asynchronous
reorganization of archival memory outside the main interaction loop.
Letta's CoALA mapping spans all three memory types (working, episodic,
semantic), but its retrieval is Naive RAG---agent-driven tool calls
with single-signal vector similarity. The OS metaphor proved
influential: MemoryOS~\cite{kang2025memoryos} and OpenViking both
adopted tiered hierarchies, though with different management
strategies.

\paragraph{Hindsight.}
Hindsight~\cite{latimer2025hindsight} introduces the most novel memory
decomposition in the production landscape: four logical networks for
\emph{world} (objective facts), \emph{experience} (agent
interactions), \emph{opinion} (subjective beliefs with confidence
scores), and \emph{observation} (preference-neutral entity summaries).
The opinion network is architecturally unique---each belief carries an
explicit confidence score that evolves as new evidence arrives,
enabling the system to distinguish between well-established facts and
tentative assessments.

Retrieval uses two specialized components: Tempr (temporal entity
memory priming) for time-aware recall and Cara (coherent adaptive
reasoning) for multi-step inference across networks. Hindsight does
not implement explicit forgetting; instead, belief confidence scores
decay or strengthen as evidence accumulates, producing an emergent
form of selective retention. The authors claim state-of-the-art
results: 91.4\% on LongMemEval and 89.61\% on LoCoMo. The four-network
decomposition is finer-grained than most systems' two- or three-store
architectures, but Hindsight lacks a graph overlay and does not
implement principled forgetting beyond confidence evolution.

\paragraph{ChatGPT Memory (OpenAI).}
OpenAI's built-in memory for ChatGPT~\cite{openai2024memoryfaq}
represents the largest-scale deployment of persistent agent memory,
available to hundreds of millions of users since early 2024. The
architecture implements a two-tier system: \emph{saved memories}
(explicit facts the user asks ChatGPT to remember, persisted
indefinitely until deleted) and \emph{reference chat history} (implicit
recall from past conversations, where details change over time as the
model decides what is most helpful to retain). Saved memories function
as auto-managed custom instructions---the model can create, update,
combine, or remove them without user intervention. When both tiers are
active, relevant information from both is injected as context into new
conversations.

Notably, ChatGPT Plus and Pro users have access to automatic memory
management that prioritizes memories based on recency and frequency of
mention, deprioritizing less-relevant memories to the
``background''---effectively a rudimentary forgetting mechanism based on
single-curve decay. The system also enforces a capacity limit (a
``memory full'' state), creating implicit pressure toward concise,
high-value memories. However, there is no structured retrieval:
memories are injected as context without ranking, no graph structure
links related memories, and no consolidation promotes episodic
interactions into semantic knowledge. Mem0 reports a 26\% accuracy
improvement over OpenAI Memory on LoCoMo~\cite{choudhary2025mem0},
consistent with our finding that full-context injection degrades as
memory stores grow. In our taxonomy, ChatGPT Memory is a Naive RAG
system with implicit preference storage and rudimentary forgetting, but
no graph, no consolidation, and no hybrid retrieval.

\paragraph{Claude Memory (Anthropic).}
Anthropic provides two distinct memory mechanisms. The \emph{Claude API
Memory Tool}~\cite{anthropic2026memorytool} enables developers to give
Claude persistent memory across conversations through a file-based
storage directory. Claude autonomously decides when to save, update, or
delete memories, storing them as structured Markdown documents organized
by topic. The tool exposes three operations---create, update, and
delete---and memories are injected as context at the start of each
conversation. This is architecturally simpler than Mem0 or Zep: no
vector retrieval, no graph traversal, no scoring---every saved memory
is included in full.

At the product level, \emph{Claude Code}'s auto-memory
system~\cite{gurgone2026claudememory} maintains a
\texttt{MEMORY.md} file per project
(\texttt{\~{}/.claude/projects/<path>/memory/MEMORY.md}), injected into
the system prompt at session start. A hard 200-line limit constrains
the file; beyond this, only the first 200 lines are loaded with a
warning nudging users toward a hub-and-spoke pattern where
\texttt{MEMORY.md} serves as an index and details live in separate
topic files (this topic-file loading is behind a feature flag and not
yet active by default). Custom agents have a three-scope memory
system: project-level (git-committable, team-shared), local
(machine-only), and user-level (global across projects). Both the API
tool and Claude Code's system are canonical examples of the
MEMORY.md problem analyzed in \S\ref{sec:memorymd}: full-context
injection with no retrieval, no lifecycle management, and a linear
scaling cost in memory size.

\paragraph{Gemini Memory (Google).}
Google's Gemini~\cite{gemini2026memory} deploys persistent memory across
its consumer assistant and Workspace applications (Gmail, Docs, Sheets,
Slides, Drive).  The architecture uses vector embeddings to store
``important user facts'' extracted from conversations, injected as
context in subsequent sessions.  A ``Recent chats'' panel (rolled out
March 2026) enables users to revisit and continue prior conversations
with full context preservation across Workspace apps.  Notably, Gemini
adopts an \emph{opt-out} privacy model---memory is enabled by default,
the opposite of ChatGPT's opt-in approach.  Memory management is
currently all-or-nothing: users can disable memory entirely or clear
all memories, but selective deletion of individual memories was not
yet available during testing.  Despite Gemini's industry-leading 1M-token
context window, the memory system remains Naive RAG: stored facts are
injected without ranking, graph structure, or lifecycle operations.

\paragraph{Grok Memory (xAI).}
xAI introduced memory for Grok in April 2025~\cite{grok2025memory},
enabling the chatbot to retain user preferences, names, and topical
interests across conversations on X (formerly Twitter), Grok.com, and
mobile apps.  The implementation follows the ChatGPT pattern: the model
decides what to remember, stored memories are transparent and
individually deletable, and a ``Personalize with Memories'' toggle
provides user control.  Architecturally, Grok Memory is Naive RAG with
implicit preference storage---no retrieval ranking, no graph, no
forgetting mechanism.  Its primary distinguishing feature is integration
with the X social platform, providing conversational context that other
assistants lack.

\paragraph{Meta AI Memory.}
Meta deploys memory across WhatsApp, Instagram, Messenger, and
meta.ai~\cite{meta2025aimemory}, potentially reaching the largest user
base of any platform memory system.  Memory stores names, preferences,
and recent topics across 1:1 chats, injected at runtime into the
model's context window.  Crucially, memory content \emph{consumes}
active context tokens---unlike systems that maintain a separate memory
store, Meta AI's injected memories reduce the available window for
conversation, creating implicit pressure toward concise memory as
histories grow.  The system provides no retrieval: all stored memories
are injected regardless of relevance, placing it firmly in the Naive
RAG category.

\paragraph{Perplexity Memory.}
Perplexity~\cite{perplexity2025memory} stands apart from the other
platform memory systems with two distinctive features.  First,
\emph{cross-model portability}: memory persists across every model
Perplexity offers (Claude, GPT, Gemini, and others), allowing users
to switch models between questions while retaining their full context
history---unique among all platforms surveyed.  Second,
\emph{citation-based transparency}: when memory influences a response,
Perplexity explicitly shows which memories contributed alongside
traditional web sources, making memory's impact auditable.  The
underlying architecture retrieves relevant context from a user's memory
store rather than injecting all memories, placing Perplexity in the
Advanced RAG category---the only consumer platform to exceed Naive RAG
for memory retrieval.  Like its peers, Perplexity lacks graph structure,
consolidation, and principled forgetting.

\paragraph{memU (NevaMind).}
memU~\cite{memu2026} is an agentic memory framework designed for 24/7
proactive agents.  Its distinguishing contribution is a hierarchical
file-system metaphor: raw interactions are transformed into structured,
searchable intelligence organized as a virtual file system, supporting
both embedding-based (RAG) and non-embedding (LLM) retrieval.  memU
receives multimodal inputs (conversations, documents, images) and
organizes them into daily notes, long-term memory, and extracted user
profiles.  A fourth ``intention layer'' (in development) analyzes
accumulated memories to predict user needs before they are expressed,
representing the most ambitious proactive memory capability in the
surveyed systems.  The framework claims approximately 10$\times$ token
reduction versus comparable usage patterns.  memU Bot extends the
framework into an enterprise-ready OpenClaw alternative with persistent
cross-session memory and team-scale deployment.

\paragraph{ReMe (MemoryScope).}
ReMe~\cite{reme2025} (formerly MemoryScope, ModelScope) provides a dual
memory architecture distinguishing \emph{personal memory} (user
preferences and traits) from \emph{task memory} (learned strategies from
prior executions).  Task memory implements four patterns: success
recognition, failure analysis, comparative learning across sampling
trajectories, and validation of extracted memories.  The framework
supports file-based and vector-based storage backends (Elasticsearch,
ChromaDB) with hybrid retrieval (vector + BM25).  ReMe integrates into
the AgentScope runtime as its default long-term memory service.  The
dual personal--task decomposition is architecturally distinct from most
systems' episodic--semantic split, treating memory \emph{function} rather
than memory \emph{type} as the primary organizing principle.

\paragraph{MIRIX.}
MIRIX~\cite{wang2025mirix} is a modular, multi-agent memory system with
six distinct memory types: Core, Episodic, Semantic, Procedural, Resource
Memory, and Knowledge Vault.  Its distinguishing contribution is
multimodal memory: unlike prior approaches confined to text, MIRIX
captures rich visual and multimodal experiences with temporally grounded
interactions.  A multi-agent framework dynamically coordinates updates
and retrieval across memory types.  On ScreenshotVQA, MIRIX achieves
35\% higher accuracy than the RAG baseline while reducing storage by
99.9\%; on LoCoMo, it attains 85.4\%, surpassing existing baselines.

%% ===============================================================
\subsection{Standalone Memory Servers}
\label{sec:standalone}

Several standalone servers provide memory as a network service,
typically via REST or MCP endpoints, without being tied to a specific
agent framework. \textbf{Motorhead} (Metal, Rust) implements a
minimal conversation buffer with incremental summarization on Redis:
when a configurable window (default 12 messages) fills, the oldest
half is summarized and the summary is incrementally updated.
\textbf{Engram} (Go) takes an agent-trusting approach---the LLM
client decides what to remember, storing curated summaries in SQLite
with FTS5 full-text search and no vector retrieval. Both represent
Naive RAG with no graph, no forgetting model, and no consolidation.
\textbf{OpenViking} (ByteDance/Volcengine, open-sourced January 2026)
adopts a virtual filesystem metaphor, mapping all agent context to a
\texttt{viking://} protocol with three tiers: L0 (immediate context),
L1 (session-level), and L2 (persistent archival). Built on VikingDB
infrastructure, OpenViking replaces fragmented vector storage with
hierarchical directory-based retrieval, but provides no forgetting,
no graph dynamics, and no preference learning.

\paragraph{MCP Memory Service.}
MCP Memory Service~\cite{doobidoo2026mcp} provides persistent memory
for AI agent pipelines (LangGraph, CrewAI, AutoGen, Claude) via a REST
API with knowledge graph and autonomous consolidation.  Storage uses
SQLite-vec with lightweight ONNX embeddings, achieving 5ms retrieval
latency.  The system supports over 20 AI client applications, making it
the most broadly compatible standalone memory server surveyed.  Its
\emph{Natural Memory Triggers} use AI-based heuristics (85\%+ accuracy)
combined with rule-based Context-Provider patterns (100\% guaranteed
coverage) to determine when to store new memories---a dual-strategy
approach unique among the surveyed systems.

%% ===============================================================
\subsection{Framework-Level Memory}
\label{sec:framework}

The major agent frameworks provide memory as a built-in module, but
these integrations are typically shallow---memory as an afterthought
rather than a first-class architectural concern. \textbf{LangChain}'s
memory classes (ConversationBufferMemory, ConversationSummaryMemory,
ConversationBufferWindowMemory) offer working-memory-only storage with
truncation-based forgetting; most are now deprecated in favor of
\texttt{RunnableWithMessageHistory}. \textbf{LlamaIndex} provides
composable memory blocks (StaticMemoryBlock,
FactExtractionMemoryBlock, VectorMemoryBlock) with FIFO eviction and
a configurable token ratio (70\% chat history, 30\% long-term by
default). \textbf{LangMem SDK}, LangChain's dedicated memory engine,
adds semantic memory (extracted facts and preferences) and a
distinctive \emph{procedural} memory capability: learned patterns are
saved as updated prompt instructions, modifying the agent's behavior
without explicit retrieval. None of these frameworks implements a
graph overlay, principled forgetting, consolidation, or preference
emergence. Their retrieval ranges from direct prompt injection
(LangChain) to vector similarity via LangGraph's storage layer
(LangMem), placing them uniformly in the Naive RAG category.

Two frameworks stand apart with substantially richer memory modules.
\textbf{AgentScope}~\cite{agentscope2025} (Alibaba/ModelScope) provides
layered memory: in-memory short-term storage, Redis-backed working memory
with TTL-based expiration, LLM-driven memory compression (added January
2026), and ReMe-based long-term memory with hybrid vector+BM25 retrieval.
AgentScope is one of the few frameworks with both consolidation and
forgetting, placing it in the Advanced RAG category.
\textbf{CAMEL}~\cite{camel2025} implements composable memory blocks via
the Composite pattern: \texttt{ChatHistoryMemory} (windowed history with
configurable window size), \texttt{VectorDBMemory} (vector similarity
over past interactions), and \texttt{LongtermAgentMemory} (a combined
module that augments chat history with vector retrieval).  CAMEL provides
the most flexible memory composition in the framework category, though
it lacks graph structure, consolidation, and forgetting.

%% ===============================================================
\subsection{Coding Agent Memory}
\label{sec:coding-agents}

A distinctive cluster of memory systems has emerged around coding
agents (Claude Code, Cursor, Windsurf, Augment Code), exposed primarily
via the Model Context Protocol (MCP).  While these tools originated as
coding assistants, their deployment as general-purpose agents---
particularly through MCP-based memory extensions---means their memory
architectures must serve domain-agnostic use cases.  We retain the
``coding agent'' label to reflect their primary deployment context, while
noting that the memory systems built for them (memsearch, QMD, Basic
Memory) are architecturally indistinguishable from dedicated memory
systems.  These systems share a common pattern:
replace the default full-context injection of MEMORY.md files with
ranked retrieval, typically achieving 70--96\% token savings.
\textbf{memsearch}~\cite{zilliz2026a} (Zilliz/Milvus team) extracts
the default memory file structure and overlays BM25 plus vector
retrieval with RRF fusion and progressive top-$k$ disclosure.
\textbf{QMD}~\cite{lutke2026qmd} (Tobi L\"utke, Shopify) runs three
local GGUF models for embedding (Gemma-300M), reranking
(Qwen3-Reranker-0.6B), and query expansion (1.7B parameter model),
achieving fully offline hybrid retrieval with cross-encoder reranking
and reported 96\% token reduction~\cite{levine2026}. \textbf{Beads}
(Go, on Dolt version-controlled SQL) models a dependency-aware issue
graph with four link types, enabling traversal along dependency chains
rather than pure similarity. \textbf{Basic Memory} (Python) takes a
wiki-link approach, storing structured Markdown with explicit
Observations and Relations that form a semantic knowledge graph
navigable via \texttt{memory://} URLs. All four independently adopted
RRF over the default weighted linear fusion---a convergence we
analyze further in \Cref{sec:memorymd}.

The OpenClaw ecosystem has spawned multiple alternative agents, each with
distinct memory approaches.
\textbf{OpenFang}~\cite{openfang2026} (RightNow AI, Rust) is the most
architecturally ambitious: its dedicated \texttt{openfang-memory} crate
implements SQLite-backed storage with vector embeddings, a knowledge
graph, cross-channel canonical sessions (context persists across
Telegram, Slack, Discord, etc.), and automatic LLM-based compaction.
OpenFang ships as a single 32MB binary with 40 channel adapters, making
it the most infrastructure-light system with graph capabilities.
\textbf{IronClaw}~\cite{ironclaw2026} (NEAR AI, Rust) implements
PostgreSQL with pgvector for persistent vector search, Reciprocal Rank
Fusion combining full-text search (tsvector/tsquery) and vector
similarity (cosine distance), and WASM sandboxing for all untrusted
tools.  Its credential isolation architecture (secrets are never exposed
to the LLM, injected only at the host boundary) is unique among the
surveyed systems.
\textbf{Nanobot}~\cite{nanobot2026} (HKU, Python) takes the opposite
approach: approximately 4,000 lines of core code with a
\texttt{MemoryStore} class maintaining daily notes
(\texttt{memory/YYYY-MM-DD.md}) and long-term memory
(\texttt{MEMORY.md})---no vector retrieval, no graph, just two plain
files and grep.

Four IDE-integrated memory systems emerged in late 2025 and early 2026.
\textbf{GitHub Copilot Memory}~\cite{copilot2026} introduces a
\emph{citation-based} architecture: memories are stored with references
to specific code locations, and agents verify citations in real-time
before using a memory.  In adversarial testing, agents consistently
detected contradictions between memories and code, updated incorrect
memories, and the memory pool self-healed.  This is the only system in
our survey that implements contradiction resolution through citation
verification rather than LLM-based fact comparison.  Memories
automatically expire after 28 days and are refreshed only through
re-validation.
\textbf{Cursor} provides two memory mechanisms:
\emph{Cursor Memory}~\cite{cursor2026} (IDE-level RAG indexing of the
entire repository via \texttt{@Codebase} symbols with vector embeddings)
and \emph{Cursor Rules} (persistent project/user/team instructions in
\texttt{.cursor/rules/*.mdc} files, injected as file-based context).
\textbf{Windsurf Memory}~\cite{windsurf2026} (Cascade) auto-generates
memories during conversation when it encounters useful context, with
a user toggle for autonomous memory creation and
\texttt{.windsurfrules} for persistent project rules.
\textbf{Augment Code}~\cite{augment2026} is the most
infrastructure-heavy IDE memory system, designed for enterprise
codebases exceeding 400{,}000 files.  Its Context Engine performs
real-time repository indexing across multiple repositories, tracking
cross-service dependencies and commit history.  The \emph{Memories}
feature stores project-specific facts---architecture decisions, coding
patterns, preferred libraries---that persist across sessions and are
automatically created by the agent when it encounters reusable context.
A \emph{Memory Review} interface lets users approve, edit, or reject
memories before finalization, providing a manual curation loop unique
among IDE memory systems.  Despite this sophistication, Augment's
memories lack automated lifecycle management: there is no consolidation,
no forgetting mechanism, and no contradiction resolution---memories
persist indefinitely until manually deleted.

\textbf{episodic-memory}~\cite{vincent2025episodic} (Jesse Vincent,
Superpowers ecosystem) archives Claude Code's conversation JSONL logs
into a SQLite database with \texttt{sqlite-vec} vector embeddings
(generated locally via Transformers.js), enabling semantic search over
the complete interaction history.  Its distinguishing contribution is a
\emph{subagent delegation} pattern: a specialized Haiku subagent handles
all memory retrieval and context assembly, preventing the context bloat
that arises from injecting raw conversation transcripts into the primary
agent's context window.  The system implements a session-end learning
loop in which the agent reflects on completed work---a ``reflect''
pattern independently adopted by several other systems.  The plugin is
part of the Superpowers marketplace, one of the largest Claude Code
plugin ecosystems.

\textbf{claude-mem}~\cite{newman2025claudemem} captures tool usage and
agent actions across five lifecycle hooks
(SessionStart, UserPromptSubmit, PostToolUse, Summary, SessionEnd) and
compresses raw observations into structured knowledge using Claude's
Agent SDK---a form of automated episodic-to-semantic consolidation.
Retrieval follows a token-efficient three-layer workflow:
\emph{search} returns a compact index of matching observation IDs,
\emph{review} presents titles and metadata for filtering, and
\emph{fetch} retrieves full details only for selected observations.
This progressive disclosure pattern minimizes context consumption---the
agent reads observation titles ($\sim$100 tokens) before deciding
whether to load full content ($\sim$500 tokens each).  An experimental
\emph{Endless Mode} implements a biomimetic memory architecture for
extended sessions, automatically managing context pressure through
adaptive summarization.  With over 28{,}000 GitHub stars, claude-mem is
among the most widely deployed agent memory systems.

Despite their retrieval sophistication, no IDE-integrated memory system
implements consolidation or automated forgetting.  Copilot's 28-day
expiry is the closest approximation, but it is a fixed TTL rather than
usage-adaptive decay.  This lifecycle gap will become acute as memory
stores grow unboundedly---particularly for Augment, where enterprise
teams accumulating memories across hundreds of thousands of files face
eventual retrieval degradation without principled memory management.

%% ===============================================================
\subsection{File-Based Memory}
\label{sec:filebased}

The simplest memory architecture stores knowledge in Markdown files
that the agent reads and writes directly. The dominant instance is the
\textbf{MEMORY.md pattern} used by coding agents such as OpenClaw:
a curated long-term knowledge file (MEMORY.md) plus daily interaction
logs with temporal decay, backed by a SQLite index with FTS5 and
sqlite-vec for hybrid retrieval (70\% vector, 30\% BM25 weighted
fusion). \textbf{Claudesidian} extends this pattern into Obsidian
vaults, using PARA folder organization and daily/weekly notes with
Obsidian's wiki-link graph for navigation. File-based memory provides
zero-infrastructure deployment and human-readable persistence, but
scales poorly: MEMORY.md is injected into the system prompt in full
on every turn regardless of relevance, consuming approximately 35,600
tokens per message (93.5\% waste in multi-message
sessions)~\cite{openclaw9157}. We analyze the cost, structure, and
retrieval failure modes of file-based memory in detail in
\Cref{sec:memorymd}.

%% ===============================================================
\subsection{Research Architectures}
\label{sec:research}

\paragraph{Generative Agents.}
Park et al.~\cite{park2023generative} introduced the seminal memory
architecture for LLM agents: a chronological \emph{memory stream} of
observations, from which the agent periodically generates
\emph{reflections}---higher-order thoughts that synthesize multiple
observations into abstract insights. Retrieval scores each memory by a
weighted product of three signals: \emph{recency} (exponential decay),
\emph{importance} (LLM-rated 1--10), and \emph{relevance} (cosine
similarity to the current query). This triple-scored formula became the
foundational retrieval mechanism for agent memory, adopted or adapted
by the majority of subsequent systems.

The architecture's significance lies less in its individual components
(all relatively simple) than in their composition: the reflection
mechanism implements a form of episodic-to-semantic consolidation,
where accumulated observations are distilled into reusable knowledge.
The system uses in-memory storage (ephemeral), a flat memory stream
(no typed stores), no graph overlay, and single-curve recency
decay---limitations that subsequent work has systematically addressed.
The reflection mechanism, however, inspired consolidation pipelines
in multiple production and research systems.

\paragraph{HippoRAG.}
HippoRAG~\cite{gutierrez2024hipporag} grounds its architecture in
hippocampal indexing theory from neuroscience: the LLM acts as
neocortex (processing and generating language), a parahippocampal
region (PHR) encoder detects synonymy and entity co-reference, and an
open knowledge graph acts as hippocampus (indexing and associating
memories). Retrieval uses Personalized PageRank (PPR) on the knowledge
graph, propagating activation from query-matched nodes to
transitively connected entities. This enables multi-hop reasoning that
pure embedding similarity cannot achieve---asking ``who handles auth
permissions?'' can retrieve information about ``Alice manages the auth
team'' through entity-relationship traversal rather than lexical
overlap.

HippoRAG achieves 20\% improvement on multi-hop question answering at
10--30$\times$ lower cost than iterative retrieval
approaches. HippoRAG~2~\cite{gutierrez2025rag2memory} (ICML~2025)
extends this with deeper passage integration---full passages become
first-class graph nodes alongside phrase nodes---and more effective
online LLM use for knowledge graph construction, adding episodic
storage and hybrid fusion (PPR + dense embedding) to the original
semantic-plus-graph design.  Neither version implements forgetting,
consolidation, or preference learning; both are retrieval
architectures rather than full memory lifecycle systems. PPR on a
knowledge graph is comparable in cognitive inspiration to spreading
activation on associative graphs, though the two mechanisms differ:
PPR computes a global stationary distribution, while spreading
activation models local, time-bounded propagation~\cite{collins1975}.

\paragraph{SYNAPSE.}
SYNAPSE~\cite{jiang2026synapse} constructs a unified episodic-semantic
graph and retrieves through two biologically inspired mechanisms:
\emph{spreading activation} (activation propagates from query-matched
nodes to neighbors through weighted edges, priming associated
memories) and \emph{lateral inhibition} (strongly activated nodes
suppress weakly activated competitors, sharpening retrieval focus).
Retrieval fuses geometric embedding similarity with activation-based
graph traversal in a triple-hybrid pipeline. Temporal decay weakens
links over time.

SYNAPSE achieves +7.2 F1 over prior state of the art on LoCoMo, 23\%
improvement on multi-hop reasoning, and 95\% token reduction versus
full-context approaches. The lateral inhibition mechanism is
particularly notable: it implements a form of retrieval-induced
forgetting~\cite{anderson1994}, where retrieving one memory
suppresses competing memories sharing the same cues. This is the most
neurobiologically grounded retrieval mechanism in the surveyed
systems, though SYNAPSE does not implement consolidation, principled
forgetting (beyond temporal decay), or preference learning. The
system is in-memory Python research code, not a persistent or
embeddable library.

\paragraph{Mem-alpha.}
Mem-alpha~\cite{wang2025memalpha} takes a fundamentally different
approach to memory management: instead of hand-crafting lifecycle
heuristics, it trains memory policies via reinforcement learning.
The architecture decomposes memory into three components---core
(always in context), episodic (interaction history), and semantic
(distilled knowledge)---a decomposition nearly identical to the CoALA
framework. An RL agent learns when to store, what to consolidate,
and what to forget, using downstream question-answering accuracy as
the reward signal. Trained on 30K tokens, the learned policy
generalizes to 400K+ tokens (13$\times$ extrapolation).

The RL-learned approach represents the most significant methodological
divergence from hand-crafted cognitive models. Learned policies can
potentially discover non-obvious management strategies that human
designers would not specify, but they require training data, are
opaque (the policy network's decisions are not interpretable), and may
not transfer across domains. Mem-alpha is part of a broader 2025--2026
trend: Memory-R1, MemRL~\cite{zhang2026memrl}, and AgeMem all train
memory management via RL, collectively establishing learned memory
policies as a distinct paradigm alongside rule-based and
neuroscience-inspired approaches.

\paragraph{SimpleMem.}
SimpleMem~\cite{simplemem2026} introduces an efficient memory framework
based on semantic lossless compression, implemented as a three-stage
pipeline: (1)~Semantic Structured Compression, which distills
unstructured interactions into compact, multi-view indexed memory units;
(2)~Online Semantic Synthesis, an intra-session process that integrates
related context into unified abstract representations to eliminate
redundancy; and (3)~Intent-Aware Retrieval Planning, which infers search
intent to dynamically determine retrieval scope.  SimpleMem achieves a
26.4\% average F1 improvement over baselines while reducing
inference-time token consumption by up to 30$\times$.

\paragraph{AriadneMem.}
AriadneMem~\cite{ariadnemem2026} extends graph-based memory with two
innovations: \emph{conflict-aware coarsening} for memory updates that
detects and resolves contradictions during graph construction, and
\emph{approximate Steiner completion with bridge-node discovery} to
construct connected evidence subgraphs for retrieval.  Rather than
expanding to large contexts or running iterative planning loops,
AriadneMem retrieves a compact connected subgraph and exposes explicit
reasoning paths to the generator.  It achieves the highest average F1
across all evaluated backbones (42.57 on GPT-4o vs.\ 39.06 for SimpleMem
and 36.09 for Mem0), with the largest gains on multi-hop and temporal
reasoning subsets.

\paragraph{MemR3.}
MemR3~\cite{memr3_2025} introduces closed-loop control over memory
retrieval through two mechanisms: a router that selects among retrieve,
reflect, and answer actions, and a global evidence-gap tracker that
renders the answering process transparent.  This design departs from the
standard retrieve-then-answer pipeline, minimizing unnecessary lookups
while providing a human-readable trace of the agent's decision process.

\paragraph{PersonaRAG.}
PersonaRAG~\cite{personarag2024} incorporates user-centric agents to
adapt retrieval and generation based on real-time user data.  Four
specialized agents---engagement tracking, preference analysis, context
understanding, and feedback integration---analyze interactions from
different perspectives, achieving over 5\% accuracy improvement over
baseline RAG.

\paragraph{ID-RAG.}
ID-RAG~\cite{idrag2025} addresses identity drift in long-horizon
generative agents by grounding persona in a structured identity
knowledge graph of core beliefs, traits, and values.  During the agent's
decision loop, this graph is queried to retrieve relevant identity
context, directly informing action selection.  The separation of identity
memory from other memory types is architecturally distinctive and
addresses a failure mode (persona drift) that no other surveyed system
explicitly targets.
